\section{The limit density at fixed memory}\label{sec:bulk-applications}

Fix $\varepsilon\in(0,\frac14)$, so $a=1-2\varepsilon\in(\frac12,1)$.
Conditional on $X_1=1$, $S_k/k^a$ converges almost surely to
$\mathcal L_a$~\cite[Theorem~3.7]{bercu2018}. The fair-start limit
$\mathcal L^\circ_a$ has the law of $\xi\mathcal L_a$, where $\xi$ is a fair sign independent of
$\mathcal L_a$~\cite[(1.3) and (1.6)]{glr2026}. Gu\'erin, Laulin and Raschel
prove that both limits have positive, bounded and smooth
densities~\cite[Theorem~1.3]{glr2026}, denoted by $h_a$ and $g_a$ as before.
Gu\'erin, Laulin, Raschel and Simon prove
that $\mathcal L_a$ is unimodal~\cite[Theorem~1.1]{glrs2025}.

\subsection{Modes and inflection points}\label{sec:modes}

A probability density $h$ on $\mathbb R$ is \emph{unimodal} if for some $m$
it is nondecreasing on $(-\infty,m]$ and nonincreasing on $[m,\infty)$.
Gu\'erin, Laulin, Raschel and Simon ask for a unique local maximum~\cite{glrs2025},
so we must also exclude a flat top. The density $g_a$ mixes the
unimodal density of $\mathcal L_a$ with its reflection. Such a
mixture can have more than one peak~\cite{dj1988}. We first take a fair start.
Lemma~\ref{lem:analytic-density} excludes constant intervals,
so modes are points. We consider other biases afterward.

\begin{theorem}[at most two modes]\label{thm:two-modes}
Let $a\in(\frac12,1)$. Exactly one of the following holds.
\begin{enumerate}[(a)]
\item The density $g_a$ is nondecreasing on $(-\infty,0]$ and nonincreasing on
$[0,\infty)$.
\item There is $\kappa>0$ such that $g_a$ is nondecreasing on $[0,\kappa]$ and
nonincreasing on $[\kappa,\infty)$, and $g_a(0)<g_a(\kappa)$.
\end{enumerate}
In case (b), $g_a$ is nondecreasing on $(-\infty,-\kappa]$ and nonincreasing on
$[-\kappa,0]$, and it is not unimodal. All these monotonicities are strict.
The unique local maximum is $0$ in case (a), and the only local maxima
are $\pm \kappa$ in case (b).
In case (a), there is $d>0$ such that $g_a$ is strictly concave on
$(-d,d)$ and strictly convex on each tail. In case (b), there are
$0<d_1<\kappa<d_2$ such that $g_a$ is strictly convex on $(-d_1,d_1)$ and
on each tail beyond $\pm d_2$, and strictly concave on
$(-d_2,-d_1)$ and $(d_1,d_2)$. Thus it has exactly two inflection points
in case (a) and four in case (b).
\end{theorem}

To pass from the second differences of $P_N$ to the curvature of
$g_a$, we smooth with a Gaussian. Factoring out a positive Gaussian
leaves an exponential sum whose signs Descartes' rule controls. We let
$N\to\infty$ at a fixed Gaussian width, and only afterward let
the width tend to $0$. This order lets us use weak convergence
before recovering curvature. Appendix~\ref{app:modes} proves the lemmas used below. A real
function $v$ has \emph{values at most} $\Pi$ if
$\operatorname{sp}(v(y_1),\dots,v(y_j))\preceq\Pi$ for all $y_1<\dots<y_j$, and
\emph{increments at most} $\Pi$ if
$\operatorname{sp}(v(y_2)-v(y_1),\dots,v(y_j)-v(y_{j-1}))\preceq\Pi$ for all
such points. We put $\Pi_4=(+,-,+,-)$.

\begin{proof}
Fix $a$, so that $\varepsilon=(1-a)/2\in(0,\frac14)$.

\emph{Step 1, smoothing.} Fix $\sigma>0$ and let
$\varphi(y)=e^{-y^2/(2\sigma^2)}/(\sigma\sqrt{2\pi})$. For $N\ge3$ put
$\eta=2N^{-a}$ and $z_k=(2k-N)N^{-a}$, so that $S_N/N^a=z_{A_N}$, and let
\[
 V_N(y)=\eta^{-2}\sum_{k=1}^{N-1}D_N(k)\,\varphi(y-z_k).
\]
Since $\varphi(y-z_k)=\varphi(y)e^{-z_k^2/(2\sigma^2)}e^{yz_k/\sigma^2}$ and
the $z_k$ increase with $k$, Lemma~\ref{lem:descartes} and
Theorem~\ref{thm:five-blocks} show that $V_N$ has values at most $\Pi_5$. (If
every $D_N(k)$ vanishes, then $V_N=0$ and every pattern is empty.)

\emph{Step 2, the limit $N\to\infty$.} Let $u(y)=\E\varphi(y-\mathcal L^\circ_a)$.
The derivatives of $\varphi$ are bounded, so $u$ is smooth and
$u^{(j)}(y)=\E\varphi^{(j)}(y-\mathcal L^\circ_a)$. We show that
$V_N(y)\to u''(y)$ for every $y$. Put $P_N(k)=0$ for $k\notin\{0,\dots,N\}$.
Then $D_N(k)=0$ unless $-1\le k\le N+1$, and summation by parts gives
\[
 \sum_{k=-1}^{N+1}D_N(k)\,\varphi(y-z_k)=\sum_{k=0}^NP_N(k)\,\Delta\varphi(y-z_k),
 \qquad \Delta\varphi(t)=\varphi(t+\eta)-2\varphi(t)+\varphi(t-\eta).
\]
Taylor's formula gives
$\lvert\eta^{-2}\Delta\varphi(t)-\varphi''(t)\rvert\le\eta^2\lVert\varphi''''\rVert_\infty/12$.
By symmetry, the four boundary terms on the left, with $k\in\{-1,0,N,N+1\}$,
have total absolute value at most
$\varphi(0)\,(6P_N(0)+2P_N(1))$. Here
$P_N(0)=\tfrac12(1-\varepsilon)^{N-1}$. The boundary recurrence gives
$\xi_{n+1}=(1-a/((1-\varepsilon)n))\xi_n+\varepsilon/(1-\varepsilon)$,
where $\xi_n=P_n(1)/P_n(0)$ and $\xi_1=1$.
Thus $\xi_n\le1+(n-1)\varepsilon/(1-\varepsilon)$ and
$P_N(1)=O_\varepsilon(N(1-\varepsilon)^N)$. Since $\eta^{-2}=N^{2a}/4$
grows only polynomially, the boundary terms scaled by $\eta^{-2}$
tend to $0$. Hence
\[
 V_N(y)=\E\varphi''\bigl(y-S_N/N^a\bigr)+o(1)\longrightarrow
 \E\varphi''(y-\mathcal L^\circ_a)=u''(y),
\]
because $S_N/N^a\to\mathcal L^\circ_a$ in law and $\varphi''$ is bounded and
continuous. By Lemma~\ref{lem:sp-limit}, $u''$ has values at most $\Pi_5$.

\emph{Step 3, integrating the sign bound.} By dominated convergence, $u'(y)\to0$ as
$y\to\pm\infty$. So $u$ has increments at most $\Pi_4$, by
Lemma~\ref{lem:sp-derivative}.

\emph{Step 4, the limit $\sigma\to0$.} Write $u=u_\sigma$. As $g_a$ is bounded
and continuous, $u_\sigma(y)=\int g_a(y-t)\varphi(t)\,dt$ tends to $g_a(y)$ for
every $y$. So the increments of $u_\sigma$ tend to those of $g_a$, and
Lemma~\ref{lem:sp-limit} shows that $g_a$ has increments at most $\Pi_4$.
Lemma~\ref{lem:analytic-density} also makes $g_a''$ bounded and continuous,
so $u_\sigma''=g_a''*\varphi\to g_a''$ pointwise. Hence $g_a''$ has
values at most $\Pi_5$.

\emph{Step 5, symmetry.} Let $0\le y_1<\dots<y_j$, let
$d_i=g_a(y_{i+1})-g_a(y_i)$, and let $p$ be the pattern of the $d_i$. As $g_a$
is even, the points $-y_j<\dots<-y_1<y_1<\dots<y_j$, with $-y_1$ left out if
$y_1=0$, have the increments $-d_{j-1},\dots,-d_1$, then $0$ if $y_1>0$, then
$d_1,\dots,d_{j-1}$. Their pattern is $-p$ read backward, followed by $p$. It has
$2\lvert p\rvert$ entries and is $\preceq\Pi_4$. So $p$ is empty, $(+)$,
$(-)$ or $(+,-)$. In particular there are no $0\le y<y'<y''$ with
$g_a(y)>g_a(y')<g_a(y'')$.

If every increment on $[0,\infty)$ is nonpositive, case (a)
holds. Otherwise choose $y_1<y_2$ with $g_a(y_1)<g_a(y_2)$.
Integrability supplies $y_3<y_4$, both beyond $y_2$, with
$g_a(y_3)>g_a(y_4)$. The sign bound excludes any later increase,
so $g_a$ attains a maximum on $[0,\infty)$ at some $\kappa$.
The absence of a strict valley makes it nondecreasing before
$\kappa$ and nonincreasing afterward. The positive increment forces
$\kappa>0$. If $g_a(0)=g_a(\kappa)$,
it would be constant on $[0,\kappa]$, contrary to analyticity
and the positive increment. Thus $\kappa>0$ and case (b) holds.
Evenness gives the left-half shape and excludes unimodality in (b).

Equality at two points of a monotone interval would force a constant
interval, contrary to Lemma~\ref{lem:analytic-density}. Thus the
monotonicities are strict and there is no additional local maximum.

\emph{Step 6, curvature.} The Fourier bound in
Lemma~\ref{lem:analytic-density} gives $g_a(y),g_a'(y)\to0$ as
$y\to\infty$. Evenness and the $\Pi_5$ bound give at most three sign
blocks for $g_a''$ on $(0,\infty)$. Its last sign must be positive. Otherwise $g_a'$ would eventually
be nonincreasing with limit zero, making $g_a$ eventually nondecreasing.
This contradicts $g_a>0$ and $g_a\to0$.
Analyticity and strict monotonicity imply that the first nonconstant
term in the even Taylor series of $g_a$ at zero is negative in case (a)
and positive in case (b). The same signs therefore hold for $g_a''$
near zero. In case (b), $g_a''$ must also be negative somewhere, since
$g_a'$ is positive near zero and negative somewhere beyond $\kappa$.
The positive-half-line sign patterns are consequently $(-,+)$ in
case (a) and $(+,-,+)$ in case (b). Analyticity excludes an interval
of zero curvature and makes the resulting convexities strict. In
case (b), $g_a'(d_1)>0>g_a'(d_2)$, so $d_1<\kappa<d_2$.
\end{proof}

\begin{theorem}[two modes for memory near one]\label{thm:two-modes-near-one}
If $\frac45\le a<1$, that is $0<\varepsilon\le\frac1{10}$, then case (b)
of Theorem~\ref{thm:two-modes} holds. So $g_a$ has exactly two modes.
\end{theorem}

\begin{proof}
In case (a), Khinchine's representation~\cite{dj1988} gives
$\mathcal L^\circ_a\overset{d}=UZ$, with $U$ uniform on $(0,1)$ and
independent of $Z$. Hence Cauchy--Schwarz gives
$\E(\mathcal L^\circ_a)^4\ge\frac95[\E(\mathcal L^\circ_a)^2]^2$.
The moment recursion of~\cite[Theorem~1.4]{glr2026} gives the actual ratio
\[
 \frac{\E\mathcal L_a^4}{(\E\mathcal L_a^2)^2}
 =\frac{6(2a^2+2a-1)}{4a-1}\,\operatorname{Beta}(2a,2a),
\]
where $\operatorname{Beta}$ is the beta function. For $a\ge\frac34$, its integral and
$t(1-t)\le\frac14$ give
$\operatorname{Beta}(2a,2a)\le4^{-2a+3/2}\operatorname{Beta}(\frac32,\frac32)=\pi2^{-4a}$.
Put $R(a)=(2a^2+2a-1)/(4a-1)$. On $[\frac34,1]$ we have
$(\log R)'\le2<4\log2$. The first inequality reduces to
$4a(4a^2+a-2)\ge0$. Thus $6\pi R(a)2^{-4a}$ decreases there, and for
$a\ge\frac45$ the displayed moment ratio is at most
\[
 \frac{282\pi}{55\,2^{16/5}}
 <\frac{141}{80}<\frac95,
\]
using $\pi<22/7$ and $2^{1/5}>8/7$. Even moments are unchanged by
symmetrization, so this contradicts case (a). Theorem~\ref{thm:two-modes}
then gives case (b).
\end{proof}

The moment argument gives an independent check near one.
To find the complete transition, we compare central curvature
across memories.

\begin{theorem}[the unique transition]\label{thm:threshold}
There is exactly one $a_*\in(1/2,1)$ with $g_{a_*}''(0)=0$.
It lies in $(2/3,69/100)$, and
\[
 g_a''(0)<0\quad(1/2<a<a_*),\qquad
 g_a''(0)>0\quad(a_*<a<1).
\]
Moreover, $\partial_a g_a''(0)|_{a=a_*}>0$.
The density $g_a$ is unimodal for $a\le a_*$ and bimodal for
$a>a_*$. In particular, the critical density is unimodal.
\end{theorem}

To rule out a second crossing, Appendix~\ref{app:analytic-threshold}
expresses the curvature sign as a comparison of positive weighted
means, without numerical input. On the short interval where the
sign can change, the rescaled weights have strictly increasing
likelihood ratios as memory increases. Together with the endpoint
signs, this gives exactly one crossing.
Appendix~\ref{app:threshold} separately locates it near $0.6688$.

\begin{remark}[finite-length thresholds]\label{rem:modes-threshold}
For the exact finite law, $D_N(N/2)$ changes sign at
$a\approx0.6563$ for $N=4000$ and at $a\approx0.6592$ for $N=8000$.
These computed examples are not used in any proof. They show how
slowly the finite law approaches the limiting transition.
\end{remark}

\subsection{Arbitrary initial bias}

Write $h_a$ for the density of $\mathcal L_a$ and
$g_{a,q}(x)=q h_a(x)+(1-q)h_a(-x)$ for initial bias
$\Prob(X_1=1)=q$. Thus $g_{a,1/2}=g_a$.

\begin{theorem}[the interval of bimodal biases]\label{thm:biased-limit-modes}
For every $a\in(1/2,1)$ and $q\in[0,1]$, the density $g_{a,q}$ has at
most two local maxima. There is $b_a\in[0,1/2]$ such that it has exactly
two local maxima if and only if
\[
                 b_a<q<1-b_a.
\]
Outside this open interval, including its endpoints, the density is
unimodal with a unique local maximum. The interval is empty exactly when
$a\le a_*$, in which case $b_a=1/2$.
At finite time, for every $0<\varepsilon<1$, $q\in[0,1]$ and
$N\ge1$, the count law has at most two modal plateaus, including its
endpoints.
\end{theorem}

\begin{proof}
\emph{The mode bound.} Let $P_N$ denote the count law with first-step bias $q$.
Proposition~\ref{prop:all-bias-modes} gives the finite assertion
and excludes two successive strict valleys. If $g_{a,q}$ had two,
weak convergence of $S_N/N^a$ and
averages over five short intervals would give five finite-law atoms
with those same forbidden inequalities
(Lemma~\ref{lem:weak-shape-transfer}).
Lemma~\ref{lem:analytic-density} excludes constant intervals
for every bias. Hence three local maxima $x<y<z$
would give
$x<u<y<v<z$ with
\[
 g_{a,q}(u)<\min\{g_{a,q}(x),g_{a,q}(y)\},\qquad
 g_{a,q}(v)<\min\{g_{a,q}(y),g_{a,q}(z)\}.
\]
For each pair, take a one-sided neighborhood of the smaller maximum
toward the other one. Since the density is not constant on an
interval, some point there has strictly smaller value. This gives
the two forbidden valleys.

\emph{The bias interval.} Fix $a$, write $f_q=g_{a,q}$ and $h=h_a$,
and call $x<y<z$ a strict valley if $f_q(y)<\min\{f_q(x),f_q(z)\}$.
Proposition~\ref{prop:dominant-half} and the same weak-transfer lemma show that,
for $q\ge1/2$, $f_q$ has no strict valley on $[0,\infty)$ and
$f_q(-x)\le f_q(x)$ for $x\ge0$.
Thus it has a strict valley if and only if
$f_q(x)>f_q(y)$ for some $x<y\le0$. Reflection dominance turns
such a descent into the valley $(x,y,-x)$. Conversely, a valley
must start at $x<0$. If its middle point $y$ is positive,
the no-valley property at $(0,y,z)$ gives
$f_q(0)\le f_q(y)<f_q(x)$.

At $q=1$, conditional unimodality~\cite{glrs2025} excludes a valley,
so $h$ is nondecreasing on $(-\infty,0]$. For fixed $x<y\le0$, the
difference $f_q(x)-f_q(y)$ is affine in $q$ and nonpositive at $q=1$.
If positive at $q\ge1/2$, it remains positive for every $r\in[1/2,q]$.
Thus, within $[1/2,1]$, a valley persists when the bias decreases.
These biases form a relatively open set because the strict
inequalities at three fixed points survive a small change of bias.
The set excludes $1$, and reflection gives an open interval
$(b_a,1-b_a)$, possibly empty.

A strict valley gives two local maxima. Integrability and positivity
let us choose an outer point on each side below the corresponding
high point. We then maximize on the two resulting compact
intervals. Conversely, two local
maxima give a strict valley because there is no flat interval.
The mode bound excludes a third maximum. A continuous integrable
density without a strict valley is unimodal, and absence of flat
intervals makes its maximum unique. Finally the bias interval is
nonempty exactly when it contains $1/2$, so
Theorem~\ref{thm:threshold} identifies the empty case.
\end{proof}

This proves the initial-bias conjecture
of~\cite[Remark~2.1(b)]{glrs2025}, with its half-width $q(a)=1/2-b_a$.
The stronger local limit of Trinas and Valle~\cite[Theorem~3.1]{tv2026},
in an unrefereed arXiv preprint, would also transfer the finite
inequalities, but weak convergence suffices here.

\subsection{The bias boundary}

To find the endpoint, we compare the conditional density's slopes
at reflected points. Their ratio tells us how much weight the
reflected density needs to create a decrease.

\begin{proposition}[a variational equation for the bias boundary]
\label{prop:bias-boundary}
For $a\in(1/2,1)$ put $U_a=\{t>0:h_a'(t)>0\}$. With the infimum
over an empty set interpreted as $+\infty$,
\begin{equation}\label{eq:bias-boundary}
 b_a=\min\left\{\frac12,\quad
       \inf_{t\in U_a}\frac{h_a'(-t)}{h_a'(t)+h_a'(-t)}\right\}.
\end{equation}
In particular $b_a\to0$ as $a\uparrow1$. If the infimum is attained
at $t>0$ with $h_a'(t)>0$ and $b_a<1/2$, then, with $q=1-b_a$,
\[
 g_{a,q}'(-t)=g_{a,q}''(-t)=0.
\]
At this point the smaller peak disappears.
\end{proposition}

\begin{proof}
The proof of Theorem~\ref{thm:biased-limit-modes} shows that
$h_a'(-t)\ge0$ for $t\ge0$. For $q\ge1/2$ bimodality is equivalent
to a decrease somewhere on the negative half-line, hence to
$g_{a,q}'(-t)<0$ for some $t>0$. Here
\[
 g_{a,q}'(-t)=q h_a'(-t)-(1-q)h_a'(t).
\]
This expression can be negative only if $t\in U_a$, in which case
it is negative exactly when
$1-q>h_a'(-t)/(h_a'(t)+h_a'(-t))$. Taking the union over $t$ gives
\eqref{eq:bias-boundary}, with equality excluded as in the interval
theorem. At an interior minimizer the ratio has derivative zero.
Differentiating its numerator and denominator gives the two displayed
equations. The moment estimate in
Corollary~\ref{cor:bias-boundary-rate} gives $b_a\to0$.
\end{proof}

\begin{theorem}[regularity of the bias boundary]\label{thm:bias-continuity}
The function $a\mapsto b_a$ is continuous on $(1/2,1)$, is continuously
differentiable on $(a_*,1)$, and extends continuously to $b_1=0$.
It is strictly decreasing on some interval immediately above $a_*$.
For every $a>a_*$, we have $0<b_a<1/2$,
and the infimum in
\eqref{eq:bias-boundary} has a unique minimizer $t_a\in(0,\kappa_a)$.
This minimizer depends continuously on $a$ throughout $(a_*,1)$.
\end{theorem}

\begin{proof}
For $a>a_*$, the curvature partition in Theorem~\ref{thm:two-modes}
gives $g_a'(t)>0$ on $(0,\kappa_a)$ and $g_a'(t)<0$ on
$(\kappa_a,\infty)$. For $t\ge\kappa_a$ in $U_a$, the ratio in
\eqref{eq:bias-boundary} is therefore at least $1/2$.
Since $b_a<1/2$, it is enough to search $(0,\kappa_a)$
(we still have to show the infimum is attained).
Put
\[
 D_a(t)=h_a'(t)+h_a'(-t),\qquad
 R_a(t)=\frac{g_a'(t)}{D_a(t)}\quad(0<t<\kappa_a).
\]
Lemmas~\ref{lem:balanced-start} and~\ref{lem:positive-conditional-mode}
give $h_a'(-t)>0$ for every $t\ge0$.
We therefore have $D_a>0$, $0<R_a<1/2$, and
$1/2-b_a=\sup_{0<t<\kappa_a}R_a(t)$. Also $R_a(0+)=0$.

We first rule out a strict valley in the ratio. We claim that
\begin{equation}\label{eq:bias-quasiconcave}
 R_a(v)\ge\min\{R_a(u),R_a(w)\}
 \qquad(0<u<v<w<\kappa_a).
\end{equation}
Otherwise, choose $s\in(0,1/2)$ strictly between the middle value
and both outer values. With $q=1/2+s$,
$g_{a,q}'(-t)=D_a(t)(s-R_a(t))$ then has signs $-,+,-$
at $-w,-v,-u$, and is positive sufficiently close to zero on the left.
Positivity of the density
and its zero limits at infinity supply a positive derivative farther
left and a negative derivative farther right. There would be three
local maxima, contrary to Theorem~\ref{thm:biased-limit-modes}.

At $\kappa_a$, we have $D_a(\kappa_a)=2h_a'(-\kappa_a)>0$.
Thus $R_a$ extends continuously to $[0,\kappa_a]$, with value
zero at both endpoints. Its positive maximum is attained in the
interior. At that point $h_a'(-t)>0$, so the maximum is strictly
below $1/2$. Hence $b_a>0$. We need attainment here (pointwise strict
inequality on an open interval would not suffice).
Equation~\eqref{eq:bias-quasiconcave} and analyticity make this
maximum unique, since two maximizers would force a constant interval.
They also make $R_a$ strictly increasing before its maximizer and
strictly decreasing after it.

Lemma~\ref{lem:analytic-density} gives joint continuity of all the
derivatives used here. The signs of $g_a'$ on either side of its
unique positive zero give continuity of $\kappa_a$. Fix $a_0>a_*$
and choose
$0<l<t_{a_0}<r<\kappa_{a_0}$. The strict inequalities
$R_a(l),R_a(r)<R_a(t_{a_0})$ persist near $a_0$.
For $v<l$, apply \eqref{eq:bias-quasiconcave} to
$v,l,t_{a_0}$. Since $R_a(l)<R_a(t_{a_0})$, it gives
$R_a(v)\le R_a(l)$. Applying it to $t_{a_0},r,v$ gives
$R_a(v)\le R_a(r)$ for $v>r$.
Thus every nearby global maximum lies in the fixed interval $[l,r]$.
The ratios converge uniformly on this fixed compact interval,
so their maxima converge. Uniqueness gives convergence of the
maximizers too. This proves continuity above $a_*$.

At $a_*$, the curvature partition in Theorem~\ref{thm:two-modes}(a)
makes $g_{a_*}'$ strictly decrease from zero and then strictly increase
to zero. Thus $g_{a_*}'(r)<0$ for every $r>0$.
Joint continuity gives $g_a'(r)<0$ for all $a>a_*$ sufficiently close
to $a_*$, so $\kappa_a<r$. Since $r$ was arbitrary, $\kappa_a\to0$.
On the shrinking interval $[0,\kappa_a]$, joint continuity now gives
\[
 \sup_{0\le t\le\kappa_a}|g_a'(t)|\longrightarrow0,\qquad
 \inf_{0\le t\le\kappa_a}D_a(t)\longrightarrow2h_{a_*}'(0)>0,
\]
where the strict inequality is Lemma~\ref{lem:positive-conditional-mode}.
Consequently $1/2-b_a=\max R_a\to0$. Below the threshold $b_a=1/2$,
so this proves continuity at $a_*$ without the critical expansion.
Proposition~\ref{prop:bias-boundary} gives the limit at one.

For differentiability, the same three-point argument puts all nearby
maximizers in one compact interval. If $M(a)=1/2-b_a$, maximality gives,
for $s>0$,
\[
 R_{a+s}(t_a)-R_a(t_a)
 \le M(a+s)-M(a)
 \le R_{a+s}(t_{a+s})-R_a(t_{a+s}).
\]
Divide by $s$ and use the continuity of $t_a$ and
$\partial_aR_a$ from Lemma~\ref{lem:analytic-density}. The argument
for $s<0$ is the same. Thus
\begin{equation}\label{eq:bias-envelope-derivative}
 -b_a'=
 \frac{\partial_a h_a'(t_a)\,h_a'(-t_a)
       -h_a'(t_a)\,\partial_a h_a'(-t_a)}
      {\{h_a'(t_a)+h_a'(-t_a)\}^2}.
\end{equation}
The right side is continuous. No derivative of $t_a$ is needed.

Finally, put $C=\partial_a g_a''(0)|_{a=a_*}>0$ and
$H=h_{a_*}'(0)>0$, by Theorem~\ref{thm:threshold} and
Lemma~\ref{lem:positive-conditional-mode}.
Extend $R_a=g_a'/D_a$ to a fixed neighborhood of zero, where
$D_a>0$ for $a$ close to $a_*$. Evenness and
Lemma~\ref{lem:analytic-density} give, jointly as $(a,t)\to(a_*,0)$,
\[
 g_a'(t)=o(t),\qquad \partial_a g_a'(t)=(C+o(1))t,\qquad
 D_a(t)=2H+o(1),\qquad \partial_aD_a(t)=O(1).
\]
For the first two relations, integrate $g_a''$ and $\partial_a g_a''$
from zero to $t$, using $g_{a_*}''(0)=0$.
It follows that
\[
 \partial_aR_a(t)=\left(\frac{C}{2H}+o(1)\right)t>0
\]
for all sufficiently small $t>0$ and $a$ sufficiently close to $a_*$.
Since $0<t_a<\kappa_a\to0$, the envelope formula gives $b_a'<0$
on some interval immediately above the threshold.
\end{proof}

\begin{corollary}[a rate for the bias boundary near one]\label{cor:bias-boundary-rate}
Put
\[
 \rho_a=\frac{\Gamma(1+a)^2}{(2a-1)\Gamma(2a)}-1.
\]
If $\rho_a<1/12$, then
\[
 0\le b_a\le\frac{4\rho_a}{1-4\rho_a}.
\]
In particular, as $a\uparrow1$,
$b_a\le8(1-a)+O((1-a)^2)$.
\end{corollary}

\begin{proof}
The conditional limit has mean $\mu_a=1/\Gamma(1+a)>0$ and variance
$\rho_a\mu_a^2$, by the moment formulas of~\cite[Theorem~1.4]{glr2026}.
Consider the three adjacent intervals of length $\mu_a$ centered at
$-\mu_a$, $0$ and $\mu_a$. Chebyshev's inequality puts at least
$1-4\rho_a$ of the conditional law in the right interval and at most
$4\rho_a$ in the middle. For a start to the left, reflect these
intervals. After mixing, the left interval has mass at least
$(1-q)(1-4\rho_a)$ and the right at least $q(1-4\rho_a)$.
The middle still has mass at most $4\rho_a$.

For $4\rho_a/(1-4\rho_a)<q<1-4\rho_a/(1-4\rho_a)$, the middle
has less mass than either outer interval. All three intervals have
the same length. So, by continuity, we can choose a point in each
with the middle value strictly smaller than the other two.
This rules out unimodality. The two-mode bound in
Theorem~\ref{thm:biased-limit-modes} gives exactly two modes, and
therefore the stated bound on $b_a$. Finally, the logarithmic derivative at $a=1$
of $1+\rho_a$ is $2\psi(2)-2-2\psi(2)=-2$.
Thus $\rho_a=2(1-a)+O((1-a)^2)$, which gives the linear bound.
\end{proof}

\subsection{The density at the threshold}

How far apart are the two modes just after the transition?
We compute this from derivatives at the threshold. Write
\[
 C_* = \left.\partial_a g_a''(0)\right|_{a=a_*},\qquad
 B_*=-g_{a_*}^{(4)}(0),\qquad H_*=h_{a_*}'(0).
\]
These constants are positive by Theorem~\ref{thm:threshold} and
Lemmas~\ref{lem:analytic-quartic} and~\ref{lem:positive-conditional-mode}.
Table~\ref{tab:critical-constants} gives their numerical enclosures.
The coefficient $Q_*$ is defined in Proposition~\ref{prop:bias-cusp}.



\begin{proposition}[the critical shape]\label{prop:critical-shape}
The density at the threshold satisfies
\[
 g_{a_*}(u)=g_{a_*}(0)-\frac{B_*}{24}u^4+O(u^6)
 \qquad(u\to0).
\]
For $a>a_*$ let $\kappa_a>0$ be the positive mode of $g_a$. As $a\downarrow a_*$,
\begin{align*}
 \kappa_a&=\sqrt{\frac{6C_*}{B_*}}\sqrt{a-a_*}
                 +o\bigl(\sqrt{a-a_*}\bigr),\\
 g_a(\kappa_a)-g_a(0)&=\frac{3C_*^2}{2B_*}(a-a_*)^2
                 +o\bigl((a-a_*)^2\bigr).
\end{align*}
\end{proposition}

\begin{proof}
Lemma~\ref{lem:analytic-density} gives joint continuity of all spatial
derivatives and their first memory derivatives. Taylor remainders
through order four are therefore uniform on a neighborhood of $a_*$.
At criticality, evenness, $g_{a_*}''(0)=0$ and the analytic sign
$g_{a_*}^{(4)}(0)=-B_*<0$ from Lemma~\ref{lem:analytic-quartic}
give the first assertion.

Put $\delta=a-a_*>0$. Uniform Taylor expansion gives
\[
 \frac{g_a'(u)}u=C_*\delta+o(\delta)
                  -\frac{B_*+o(1)}6u^2+O(u^4).
\]
For $u=k\sqrt\delta$ this has the sign of $C_*-B_*k^2/6$
when $k$ is fixed away from $\sqrt{6C_*/B_*}$. The unique positive
mode lies between any two such choices with opposite signs, by
Theorem~\ref{thm:two-modes}. Squeezing gives the mode expansion.
Substitute it into
\[
 g_a(u)-g_a(0)=\tfrac12\bigl(C_*\delta+o(\delta)\bigr)u^2
                  -\tfrac1{24}\bigl(B_*+o(1)\bigr)u^4+O(u^6)
\]
to obtain the height expansion.
\end{proof}

\begin{proposition}[how the bias interval opens]\label{prop:bias-cusp}
With $C_*,B_*,H_*$ as defined above, put
\[
 Q_* = \frac{\sqrt{2}\,C_*^{3/2}}{3H_*\sqrt{B_*}},
 \qquad Q_*\approx9.867.
\]
As $a\downarrow a_*$,
\begin{equation}\label{eq:bias-cusp}
 \frac12-b_a=Q_*(a-a_*)^{3/2}+o\bigl((a-a_*)^{3/2}\bigr).
\end{equation}
The unique minimizer $t_a$ of
Theorem~\ref{thm:bias-continuity} satisfies
\[
 t_a=\sqrt{\frac{2C_*}{B_*}}\sqrt{a-a_*}
       +o\bigl(\sqrt{a-a_*}\bigr)
     =\frac{\kappa_a}{\sqrt3}+o\bigl(\sqrt{a-a_*}\bigr).
\]
At $q=1-b_a$, the stationary point is $-t_a$.
\end{proposition}

\begin{proof}
Lemma~\ref{lem:analytic-density} supplies the required joint regularity.
Theorem~\ref{thm:threshold} and
Lemma~\ref{lem:positive-conditional-mode} give $C_*>0$ and $H_*>0$.

Write $\delta=a-a_*$. By \eqref{eq:bias-boundary},
\[
 \frac12-b_a=\max\left\{0,\sup_{t\in U_a}
     \frac{g_a'(t)}{h_a'(t)+h_a'(-t)}\right\}.
\]
The fair-shape theorem makes the numerator nonpositive for $t\ge \kappa_a$.
Proposition~\ref{prop:critical-shape} gives
$\kappa_a=O(\sqrt\delta)$, and $H_*>0$ makes both conditional derivatives
positive on this shrinking interval. Uniform Taylor expansion gives,
for bounded $v\ge0$,
\[
 \frac{g_a'(v\sqrt\delta)}{h_a'(v\sqrt\delta)+h_a'(-v\sqrt\delta)}
 =\frac{\delta^{3/2}}{2H_*}
       \left(C_*v-\frac{B_*v^3}{6}\right)+o(\delta^{3/2}).
\]
The cubic is largest at the unique positive point
$v=\sqrt{2C_*/B_*}$. Maximizing on $[0,\kappa_a/\sqrt\delta]$, whose
right endpoint tends to $\sqrt{6C_*/B_*}$, proves \eqref{eq:bias-cusp}.
Since the ratio vanishes at both endpoints, its positive maximum
is attained inside $(0,\kappa_a)$. The convergence is uniform
and the limiting cubic has a unique maximum, so every scaled
maximizer tends to
$\sqrt{2C_*/B_*}$. This proves the assertion about $t_a$.


\end{proof}

The two powers come from the width and height of the same cubic.
Near zero the derivative is approximately $C_*(a-a_*)t-B_*t^3/6$.
Its positive part has width of order $\sqrt{a-a_*}$ and height of
order $(a-a_*)^{3/2}$. The width sets the distance between the modes.
The height sets how much bias removes the smaller peak.

